
\subsection{循环对称性与定向}
\label{sec:8-4-2}

本小节描述定向在循环对称性\liftIndex{I113}下的行为。首先引入相交配对
$\langle\ ,\ \rangle$ 的一种版本。

\begin{definition}{8.4.6}\label{def:8-4-6}
设 $P_1,P_2$ 是一对彼此横截、维数互补的光滑奇异单形。定义
\[
 \langle P_1,P_2\rangle:=\#(P_1\times_LP_2).
\]
\end{definition}

\begin{remark}{8.4.7}\label{rem:8-4-7}
\textup{(1)} 令 $N_i$ 为 $P_i$ 在 $L$ 中的法丛，并把 $N_i$ 延拓到
$P_i$ 的一个管状邻域上，其中 $i=1,2$。于是，$P_i$ 可写成该管状邻域上
$N_i$ 的重言截面 $s_i$ 的零点集。%
\translatornote{源书把承载 $s_i$ 的丛写成未定义的 $E_i$；本段只定义了
法丛 $N_i$，紧接的仓西邻域及直和截面也都使用 $N_i$，故改回 $N_i$。}
由约定~\ref{conv:8-2-1}(4)，有
\[
\begin{aligned}
 P_1\times_LP_2
 &=(s_1;N_1\longrightarrow L)\times_L
   (s_2;N_2\longrightarrow L)\\
 &=(-1)^{\deg P_2\cdot\deg P_1}(s_1\oplus s_2)^{-1}(0).
\end{aligned}
\]
计入约定 8.6.1，得到
\[
 \langle P_1,P_2\rangle
 =(-1)^{\deg P_1\cdot\deg P_2}\#(P_1\cap P_2).
\]

分别把 $P_1,P_2$ 看作流 $T(P_1),T(P_2)$。由于已经假设 $P_1,P_2$
横截相交，乘积 $T(P_1)\wedge T(P_2)$ 有定义。于是上述观察可重述为
\[
 \langle P_1,P_2\rangle
 =(-1)^{\deg P_1\cdot\deg P_2}
   \int_L T(P_1)\wedge T(P_2)
 =\int_L T(P_2)\wedge T(P_1).
\]

\textup{(2)} 为了在 $C^\bullet(L)$ 上定义该配对，我们像处理
$A_\infty$ 结构那样，取扰动后的相交数。
\end{remark}

\begin{proposition}{8.4.8}\label{prop:8-4-8}
置 $\deg'=\deg+1$。则
\[
\begin{aligned}
 &\bigl\langle P_0,\mathfrak m_{k,\beta}(P_1,\ldots,P_k)\bigr\rangle\\
 &\quad={}
 (-1)^{\deg' P_k(\deg' P_0+\cdots+\deg' P_{k-1})}
 \bigl\langle P_k,\mathfrak m_{k,\beta}
 (P_0,\ldots,P_{k-1})\bigr\rangle.
\end{aligned}
\]
\end{proposition}

与引理~\ref{lem:8-4-3} 一样，我们使用并非主分支的有边稳定映射
模空间；也就是说，边界标记点不必遵循边界上的典范循环次序。这里施行
交换第 $0$、第 $1$ 个标记点的置换。令
$\mathcal M'_{k+1}(\beta)$ 为
$(z_1,z_0,z_2,\ldots,z_k)$ 的次序与逆时针定向相容的模空间，并置
\[
\begin{aligned}
 &\mathcal M'_{k+1}(\beta;P_0,P_2,\ldots,P_k)\\
 &\quad={}
 (-1)^{\epsilon'}\mathcal M'_{k+1}(\beta)
 \mathop{\times}_{L\times\cdots\times L}
 (P_0\times P_2\times\cdots\times P_k),
\end{aligned}
\]
其中%
\translatornote{源书式 (8.4.9) 的内层求和写作 $\deg P_j$，使内指标
$i$ 没有出现；把定义~\ref{def:8-4-1} 应用于依次排列的
$P_0,P_2,\ldots,P_k$，以及本证明末尾对 $\epsilon+\epsilon'$ 的化简，
都唯一要求这里为 $\deg P_i$。}
\begin{equation}
 \epsilon'=(n+1)\left(
 (k-1)\deg P_0+
 \sum_{j=2}^{k-1}\sum_{i=2}^{j}\deg P_i
 \right).
 \tag{8.4.9}\label{eq:8-4-9}
\end{equation}
比较定义~\ref{def:8-4-1}。

\begin{lemma}{8.4.10}\label{lem:8-4-10}
有
\[
\begin{aligned}
 &P_0\times_L
 \mathcal M_{k+1}(\beta;P_1,P_2,\ldots,P_k)\\
 &\quad={}
 (-1)^{\deg'P_0\deg'P_1}
 P_1\times_L
 \mathcal M'_{k+1}(\beta;P_0,P_2,\ldots,P_k).
\end{aligned}
\]
\end{lemma}

\begin{proof}
用 $L_i$ 表示与 $P_i$ 对应的评价映射之靶，并写成
\[
 \mathcal M_{k+1}(\beta)
 =L_0\times{}^\circ\mathcal M_{k+1}(\beta)^\circ
  \times L_1\times L_2\times\cdots\times L_k,
\]
以及
\[
 \mathcal M'_{k+1}(\beta)
 =L_1\times{}^\circ\mathcal M'_{k+1}(\beta)^\circ
  \times L_0\times L_2\times\cdots\times L_k.
\]
注意，因为交换了前两个标记点，
\[
 \mathcal M_{k+1}(\beta)=-\mathcal M'_{k+1}(\beta),
\]
并且
\[
 {}^\circ\mathcal M'_{k+1}(\beta)^\circ
 =(-1)^\gamma{}^\circ\mathcal M_{k+1}(\beta)^\circ,
\]
其中
\[
 \gamma=(\dim L)^2+1\equiv n+1\pmod2.
\]
计入式~\eqref{eq:8-2-7}，由定义~\ref{def:8-4-1}，有
\[
\begin{aligned}
 \mathcal M_{k+1}(\beta;P_1,\ldots,P_k)
 ={}&(-1)^\epsilon L_0
 \times{}^\circ\mathcal M_{k+1}(\beta)^\circ\\[-2pt]
 &{}\times L_1\times{}^\circ P_1\times\cdots
 \times L_k\times{}^\circ P_k,
\end{aligned}
\]
其中
\[
 \epsilon=(n+1)
 \sum_{j=1}^{k-1}\sum_{i=1}^{j}\deg P_i,
\]
并且
\[
\begin{aligned}
 &\mathcal M'_{k+1}(\beta;P_0,P_2,\ldots,P_k)\\
 &\quad={}
 (-1)^{\epsilon'}L_1
 \times{}^\circ\mathcal M'_{k+1}(\beta)^\circ
 \times L_0\times{}^\circ P_0
 \times L_2\times{}^\circ P_2\times\cdots
 \times L_k\times{}^\circ P_k,
\end{aligned}
\]
这里 $\epsilon'$ 由式~\eqref{eq:8-4-9} 给出。%
\translatornote{源书在 $L_0\times{}^\circ P_0$ 后把 $L_2$ 误印成
$P_2$，从而出现两个相邻的 $P_2$ 因子；紧接的定向比较始终使用
$(L_2\times{}^\circ P_2)$，故补回 $L_2$。}
于是
\[
\begin{aligned}
 &P_0\times_L\mathcal M_{k+1}(\beta;P_1,P_2,\ldots,P_k)\\
 &={}
 (-1)^\epsilon(P_0^\circ\times L_0)
 \times{}^\circ\mathcal M_{k+1}(\beta)^\circ
 \times(L_1\times{}^\circ P_1)\\[-2pt]
 &\qquad{}\times(L_2\times{}^\circ P_2)\times\cdots
 \times(L_k\times{}^\circ P_k)\\
 &=(-1)^{\epsilon+\delta}(L_1\times{}^\circ P_1)
 \times{}^\circ\mathcal M_{k+1}(\beta)^\circ
 \times(P_0^\circ\times L_0)\\[-2pt]
 &\qquad{}\times(L_2\times{}^\circ P_2)\times\cdots
 \times(L_k\times{}^\circ P_k)\\
 &=(-1)^{\epsilon+\delta+\gamma}(P_1^\circ\times L_1)
 \times{}^\circ\mathcal M'_{k+1}(\beta)^\circ
 \times(L_0\times{}^\circ P_0)\\[-2pt]
 &\qquad{}\times(L_2\times{}^\circ P_2)\times\cdots
 \times(L_k\times{}^\circ P_k)\\
 &=(-1)^{\epsilon+\delta+\gamma+\epsilon'}
 P_1\times_L\mathcal M'_{k+1}(\beta;P_0,P_2,\ldots,P_k),
\end{aligned}
\]
其中
\[
 \delta=p_0p_1+(n+1)(p_0+p_1)k,
 \qquad p_i=\dim P_i.
\]
在第二个等号中，我们用到
\[
\begin{aligned}
 \dim{}^\circ\mathcal M_{k+1}(\beta)^\circ
 &=n+\mu(\beta)+k+1-3-n(k+1)\\
 &\equiv(n+1)k\pmod2.
\end{aligned}
\]
在第三个等号中，我们用了
\[
 L\times{}^\circ P=P=P^\circ\times L.
\]
注意
\[
 \epsilon+\epsilon'
 \equiv(n+1)(k+1)(p_0+p_1)\pmod2,
\]
故
\[
\begin{aligned}
 \epsilon+\delta+\gamma+\epsilon'
 &\equiv(n+1+p_0)(n+1+p_1)\\
 &\equiv(\deg P_0+1)(\deg P_1+1)\pmod2.
\end{aligned}
\]
引理~\ref{lem:8-4-10} 得证。
\end{proof}

\begin{proof}[命题~\ref{prop:8-4-8} 的证明]
命题~\ref{prop:8-4-8} 由引理~\ref{lem:8-4-3} 和
引理~\ref{lem:8-4-10} 立即得到。
\end{proof}
